%%%%%%%%%%%%%%%%%%%%%%%%%%%%%%%%%%%%%%%%%%%%%%%%%%%%%%%%%%%%
%% Chapter 6 — Control Framework and Simulation Integration
%% (12–18 pages)
%% Purpose: explain the complete hierarchical controller and
%% present PPO1 and PPO2 as two controlled experimental
%% pipelines rather than a collection of disconnected models.
%% Main output: reproducible controller architecture, PPO task
%% definition, integration details, fair configuration of all
%% methods.
%%%%%%%%%%%%%%%%%%%%%%%%%%%%%%%%%%%%%%%%%%%%%%%%%%%%%%%%%%%%
\chapter{Control Framework and Simulation Integration}
\label{chap:control}

\section{Nominal Task-Space Controller}
\label{sec:ctrl:dls}
% - Reference-pose generation and task-space error.
% - DLS inverse differential kinematics and damping strategy.
% - Velocity/position-increment limits; mapping to cylinder targets.
% - DLS as nominal structured controller, not the adaptive contribution.
\todo{write}

\section{Residual PPO Formulation}
\label{sec:ctrl:ppo}
% - What PPO corrects: task-space commands, joint/cylinder targets, or
%   low-level commands.
% - Observation groups, action scaling, action filtering, temporal
%   history.
% - Reward terms: tracking, smoothness, stability, limits,
%   pressure/effort penalties, residual magnitude.
% - Training curriculum, episode termination, policy network, PPO
%   hyperparameters.
\todo{write}

\section{PPO1-A: Nominal ID-FD Pipeline}
\label{sec:ctrl:ppo1a}
% - ID and FD trained on nominal data; no AAM, no latent conditioning.
% - Purpose: nominal learned-dynamics baseline under residual PPO.
\todo{write}

\section{PPO1-B: Randomized/Robust ID-FD Pipeline}
\label{sec:ctrl:ppo1b}
% - ID and FD trained/fine-tuned with randomized and informative data;
%   no AAM on the FD.
% - Purpose: does data diversity alone improve the ID-FD pipeline?
\todo{write}

\section{PPO2-A: Nominal ROM Pipeline}
\label{sec:ctrl:ppo2a}
% - ROM trained on nominal data, used with the PPO2 policy architecture.
% - No latent conditioning and no AAM.
% - Purpose: baseline for the reduced-order architecture.
\todo{write}

\section{PPO2-B: Randomized/Robust ROM Pipeline}
\label{sec:ctrl:ppo2b}
% - ROM trained across randomized systems, no online latent
%   conditioning.
% - Purpose: isolate robustness from adaptation.
\todo{write}

\section{PPO2-C: AAM-Conditioned ROM Pipeline}
\label{sec:ctrl:ppo2c}
% - AAM estimates current latent dynamics from recent observations;
%   the latent estimate conditions the ROM online.
% - Purpose: complete proposed adaptive method.
\todo{write}

\section{Optional Oracle-Latent Diagnostic}
\label{sec:ctrl:oracle}
% - True simulation latent/context supplied to the ROM as an
%   upper-bound diagnostic only (PPO2-D).
% - Gap between oracle and AAM separates ROM capacity limitations from
%   adaptation-estimation limitations.
\todo{write}

\section{IsaacLab Integration and Timing}
\label{sec:ctrl:integration}
% - Control and physics time steps, action decimation, observation
%   timing, history-buffer updates.
% - Model inference order and handling of previous actions.
% - Device placement, vectorized environments, checkpoint loading,
%   runtime cost.
\todo{write}

\section{Stability and Constraint Handling}
\label{sec:ctrl:constraints}
% - Joint, cylinder, velocity, pressure, valve limits.
% - Action clipping, rate limits, termination conditions, numerical
%   safeguards.
% - How residual authority is bounded so PPO cannot completely
%   override the nominal controller.
\todo{write}
